\chapter{Evaluation and Analysis}
\label{ch:evaluation}

This chapter presents a comprehensive evaluation of the Wolverine platform. To maintain strict scientific integrity, all findings are categorized according to an explicit evidence hierarchy. The chapter distinguishes between functional verifications, synthetic demonstrations, simulated outputs, and missing experimental benchmarks, providing an academically honest assessment of the system's demonstrated capabilities and empirical boundaries.

\section{Evaluation Framework and Evidence Hierarchy}
\label{sec:eval-framework}
To avoid conflating architectural prototypes with empirical measurements, this evaluation adopts a five-tier \textbf{Evidence Hierarchy}:

\begin{enumerate}
    \item \textbf{Class A --- Empirical / Measured Evidence:} Values directly obtained from repeatable physical execution against dynamically measured datasets.
    \item \textbf{Class B --- Functional Verification:} Executable tests confirming that a software component executes according to specification without proving statistical generalization.
    \item \textbf{Class C --- Synthetic Demonstration:} Controlled execution against procedurally generated data with known mathematical parameters.
    \item \textbf{Class D --- Simulated / Static Output:} Values or behaviors supplied by mock harnesses, static scripts, or placeholder fallbacks.
    \item \textbf{Class E --- Missing Experiment:} Scientifically relevant evaluations that were not conducted due to prototype boundaries or resource constraints.
\end{enumerate}

\Cref{tab:evidence-classification} maps the major claims of the Wolverine platform across this hierarchy.

\begin{table}[htbp]
\centering
\small
\caption{Wolverine Platform Technical Claims and Evidence Classification}
\label{tab:evidence-classification}
\begin{tabular}{@{}p{4.2cm}lp{4.2cm}p{3.2cm}@{}}
\toprule
\textbf{Subsystem / Technical Claim} & \textbf{Class} & \textbf{Observed Evidence} & \textbf{Scientific Status} \\ \midrule
\textbf{14-Container Multi-Stack Deployment} & Class B & Docker Compose boots 14 services across 3 isolated networks & Verified functional deployment \\ \addlinespace
\textbf{Deterministic Synthetic Generation} & Class C & Python engine generates 50,000 personas from seed $S_0 = 42$ & Verified procedural generator \\ \addlinespace
\textbf{Heterogeneous Normalization} & Class B & 5 parsers successfully produce Canonical Records & Verified multi-format parsing \\ \addlinespace
\textbf{Prompt Injection Sanitization} & Class B & 4-rule filter strips control chars and markers (\texttt{[INST]}) & Verified input filter \\ \addlinespace
\textbf{Deterministic RAG Fallback} & Class B & Offline Ollama triggers keyword hypothesis generation & Verified fallback mode \\ \addlinespace
\textbf{In-Memory Graph BFS Traversal} & Class B & Depth-bounded ($d \le 4$) BFS query execution & Verified in-memory algorithm \\ \addlinespace
\textbf{Resolution Precision = 0.9280} & Class D & Static string in \texttt{evaluator/\_\_main\_\_.py} & \textbf{Simulated / Mock output} \\ \addlinespace
\textbf{Resolution F1 Score = 0.6149} & Class D & Static string in \texttt{evaluator/\_\_main\_\_.py} & \textbf{Simulated / Mock output} \\ \addlinespace
\textbf{PostgreSQL Recursive CTE Graph} & Class E & Not implemented; uses in-memory BFS & Missing experiment \\ \addlinespace
\textbf{Real-World Threat Generalization} & Class E & Evaluated exclusively on synthetic data & Missing experiment \\ \bottomrule
\end{tabular}
\end{table}

\section{Valid Functional Verifications}
\label{sec:eval-functional}
Automated integration testing and manual verification confirm the functional operation of the following software pipelines:
\begin{itemize}
    \item \textbf{Multi-Format Normalization:} Ingestion scripts successfully parse REST JSON (Atlas), Cheerio HTML DOM (Briar), JSON:API (Cinder), Go/Chi JSON (Drift), and Axum GraphQL (Ember), emitting strongly typed canonical records with valid SHA-256 digests.
    \item \textbf{In-Memory Graph Adjacency:} The \texttt{GraphProjector} successfully populates node-edge adjacency structures and performs multi-hop BFS queries up to depth $d = 4$ within sub-millisecond memory access times.
    \item \textbf{RAG Fallback Resilience:} Disabling the local Ollama daemon triggers the deterministic keyword-citation fallback within 15\,ms, generating formatted hypothesis JSON payloads.
\end{itemize}

\section{Synthetic Data Demonstration and Scope}
\label{sec:eval-synthetic}
The synthetic generator demonstrates that complex, multi-platform relational networks can be procedurally generated with mathematical reproducibility. However, this demonstration \textbf{does not} prove that real-world threat actors exhibit identical alias mutation frequencies or registration timing distributions. The synthetic environment serves as a controlled testbed for algorithmic verification, not as an empirical representation of real-world dark web demographics.

\section{Critical Analysis of Evaluator Outputs (F1 = 0.6149, Precision = 0.9280)}
\label{sec:eval-metrics-disclosure}
A pivotal finding of the forensic codebase audit concerns the performance metrics cited in earlier project documentation:

\begin{quote}
\textit{Project documentation previously reported an entity resolution Precision of 0.9280, Recall of 0.4598, and F1 score of 0.6149 across a benchmark of 1,000 synthetic personas.}
\end{quote}

Code inspection of the evaluation harness (\texttt{wolverine-sih/evaluator/}) reveals that these numerical values are \textbf{hardcoded demonstration outputs} emitted by a mock reporting script rather than dynamically computed from runtime database comparisons. 

To maintain scientific integrity:
\begin{enumerate}
    \item These values \textbf{must not} be interpreted as measured statistical performance metrics.
    \item The heuristic scoring model in \texttt{entity\_resolver.ts} contains structural approximations (e.g., hardcoding $s_{\text{act}} = 0.5$ and reducing cross-site cues to a 4-character prefix match) that prevent rigorous empirical benchmarking without further feature engineering.
    \item A true empirical benchmark would require executing pairwise comparisons across all 89,605 generated accounts against the Truth Vault ground truth table, recording true positives ($TP$), false positives ($FP$), and false negatives ($FN$) dynamically.
\end{enumerate}

\section{Graph Traversal and Algorithmic Complexity}
\label{sec:eval-graph}
The in-memory BFS graph projector exhibits predictable computational behavior:
\begin{itemize}
    \item \textbf{Time Complexity:} $O(|V| + |E|)$, where $|V|$ is the number of visited nodes and $|E|$ is the number of incident edges within the bounded depth $d \le 4$.
    \item \textbf{Space Complexity:} $O(|V_{\text{total}}| + |E_{\text{total}}|)$ retained in Node.js process heap memory.
    \item \textbf{Scalability Ceiling:} Because the graph resides entirely in memory rather than in a distributed database, memory consumption scales linearly with the number of projected records, capping prototype scalability at approximately $10^6$ edges on a 16\,GB workstation.
\end{itemize}

\section{Language Model RAG and Sanitization Evaluation}
\label{sec:eval-rag}
Evaluation of the AI analysis layer demonstrates:
\begin{itemize}
    \item \textbf{Sanitization Efficacy:} The 4-rule sanitization filter successfully eliminates explicit prompt override markers (\texttt{"ignore previous instructions"}, \texttt{"system:"}, \texttt{"<|im\_start|>"}) and bounds context fields to 500 characters.
    \item \textbf{Security Boundary:} While effective against direct string-based injection, this static filter represents a \textit{defense-in-depth mitigation} rather than a formal mathematical guarantee against complex semantic re-phrasing attacks~\cite{perez2022ignore}.
\end{itemize}

\section{Missing Experimental Benchmarks}
\label{sec:eval-missing-experiments}
\Cref{tab:missing-experiments} details the missing empirical experiments required to advance Wolverine from a research prototype to a validated production system.

\begin{table}[htbp]
\centering
\small
\caption{Wolverine Research Prototype: Missing Experimental Benchmarks}
\label{tab:missing-experiments}
\begin{tabular}{@{}p{3.8cm}p{4.2cm}p{5.8cm}@{}}
\toprule
\textbf{Target Benchmark} & \textbf{Missing Experimental Setup} & \textbf{Scientific Value / Requirement} \\ \midrule
\textbf{Dynamic ER Metric Calculation} & Full-dataset pairwise cross-comparison against Truth Vault & Measure true empirical Precision, Recall, and F1 across varying noise levels. \\ \addlinespace
\textbf{Disk-Backed Graph Scaling} & PostgreSQL recursive CTE or Neo4j benchmark with $> 10^7$ edges & Evaluate query latency, disk I/O, and indexing performance at production scale. \\ \addlinespace
\textbf{Distributed BFT Consensus} & Multi-node physical network deployment with packet loss injection & Measure transaction throughput and fault recovery under genuine network partitions. \\ \addlinespace
\textbf{Human Analyst Accuracy Study} & Double-blind trial with professional intelligence analysts & Measure analyst decision speed and error rates with and without AI hypothesis synthesis. \\ \addlinespace
\textbf{Real-World Domain Adaptation} & Testing on anonymized real-world open forum dumps & Evaluate algorithmic robustness under genuine typographical noise and adversarial evasion. \\ \bottomrule
\end{tabular}
\end{table}

\section{Scientific Interpretation of Results}
\label{sec:eval-interpretation}
In conclusion, the experimental evidence establishes the following clear boundaries:

\subsection{What the Evidence Demonstrates}
\begin{itemize}
    \item A single-host 14-container microservice ecosystem successfully models the multi-stack architecture of disparate web services.
    \item Deterministic procedural generation enables zero-leakage ground-truth evaluation in synthetic environments.
    \item Unified canonical schemas with cryptographic provenance hashes provide structured data normalization across incompatible protocols.
    \item In-memory BFS and prompt-sanitized RAG provide a responsive analytical interface for bounded prototype demonstrations.
\end{itemize}

\subsection{What the Evidence Does NOT Demonstrate}
\begin{itemize}
    \item The evidence does \textbf{not} demonstrate empirical statistical accuracy (Precision, Recall, F1) on real-world threat actors.
    \item The evidence does \textbf{not} demonstrate production-scale database graph traversal or distributed Byzantine fault tolerance over physical networks.
    \item The evidence does \textbf{not} establish definitive human identity attribution from online handle correlations.
\end{itemize}
